\section{Injectives and adjoint functors}
\label{section-adjoint}

\noindent
Here are some lemmas on adjoint functors and their relationship with
injectives. See also Lemma \ref{lemma-adjoint-get-abelian}.

\begin{lemma}
\label{lemma-adjoint-preserve-injectives}
\begin{slogan}
A functor with an exact left adjoint preserves injectives
\end{slogan}
Let $\mathcal{A}$ and $\mathcal{B}$ be abelian categories.
Let $u : \mathcal{A} \to \mathcal{B}$ and
$v : \mathcal{B} \to \mathcal{A}$ be additive functors with
$u$ right adjoint to $v$. Consider the following conditions:
\begin{enumerate}
\item[(a)] $v$ transforms injective maps into injective maps,
\item[(b)] $v$ is exact, and
\item[(c)] $u$ transforms injectives into injectives.
\end{enumerate}
Then (a) $\Leftrightarrow$ (b) $\Rightarrow$ (c). If $\mathcal{A}$
has enough injectives, then all three conditions are equivalent.
\end{lemma}

\begin{proof}
Observe that $v$ is right exact as a left adjoint (Categories, Lemma
\ref{categories-lemma-exact-adjoint}). Combined with
Lemma \ref{lemma-exact-functor} this explains why (a) $\Leftrightarrow$ (b).

\medskip\noindent
Assume (a). Let $I$ be an injective object of $\mathcal{A}$.
Let $\varphi : N \to M$ be an injective map in $\mathcal{B}$ and let
$\alpha : N \to uI$ be a morphism.
By adjointness we get a morphism $\alpha : vN \to I$ and
by assumption $v\varphi : vN \to vM$ is injective.
Hence as $I$ is an injective object we get a morphism
$\beta : vM \to I$ extending $\alpha$. By adjointness
again this corresponds to a morphism $\beta : M \to uI$ extending $\alpha$.
Hence (c) is true.

\medskip\noindent
Assume $\mathcal{A}$ has enough injectives and (c) holds.
Let  $f : B \to B'$ be a monomorphism in $\mathcal{B}$, and let
$A = \Ker(v(f))$. Choose a monomorphism $g : A \to I$ with $I$ injective
(it exists by assumption). Then $g$ extends to $g' : v(B) \to I$,
whence by adjunction a morphism $B \to u(I)$. Since $u(I)$ is injective,
this morphism extends to $h : B' \to u(I)$, hence by adjunction a morphism
$k : v(B') \to I$ extending $g'$. But then $k$ "extends'' $g$,
which forces $A = 0$ since $g$ was a monomorphism. Thus (a) is true.
\end{proof}

\begin{remark}
\label{remark-need-left-exactness}
Let $R \to S$ be a ring map. Let $u : \text{Mod}_S \to \text{Mod}_R$
be $u(N) = N_R$ and let $v : \text{Mod}_R \to \text{Mod}_S$ be
$v(M) = M \otimes_R S$. Then $u$ is right adjoint to $v$, and $u$ is
exact and $v$ is right exact. But conditions (a), (b), (c) of
Lemma \ref{lemma-adjoint-preserve-injectives} do not hold in general.
For example, if $R = \mathbf{Z}$ and $S = \mathbf{Z}/p\mathbf{Z}$, then
the injective $S$-module $\mathbf{Z}/p\mathbf{Z}$ is not
an injective $\mathbf{Z}$-module. In fact, the lemma shows
all injective $S$-modules are injective as $R$-modules if and only
if $R \to S$ is a flat ring map.
\end{remark}

\begin{lemma}
\label{lemma-adjoint-enough-injectives}
Let $\mathcal{A}$ and $\mathcal{B}$ be abelian categories.
Let $u : \mathcal{A} \to \mathcal{B}$ and
$v : \mathcal{B} \to \mathcal{A}$ be additive functors.
Assume
\begin{enumerate}
\item $u$ is right adjoint to $v$,
\item $v$ transforms injective maps into injective maps,
\item $\mathcal{A}$ has enough injectives, and
\item $vB = 0$ implies $B = 0$ for any $B \in \Ob(\mathcal{B})$.
\end{enumerate}
Then $\mathcal{B}$ has enough injectives.
\end{lemma}

\begin{proof}
Pick $B \in \Ob(\mathcal{B})$.
Pick an injection $vB \to I$ for $I$
an injective object of $\mathcal{A}$.
According to Lemma \ref{lemma-adjoint-preserve-injectives}
and the assumptions the corresponding map
$B \to uI$ is the injection of $B$ into an injective object.
\end{proof}

\begin{remark}
\label{remark-faithfulness-needed}
Let $\mathcal{A}$, $\mathcal{B}$, $u : \mathcal{A} \to \mathcal{B}$ and
$v : \mathcal{B} \to \mathcal{A}$ be as
In Lemma \ref{lemma-adjoint-enough-injectives}.
In the presence of conditions (1) and (2) condition (4) is equivalent
to $v$ being faithful. Moreover, condition (4) is needed.
An example is to consider the
case where the functors $u$ and $v$ are both the zero functor.
\end{remark}

\begin{lemma}
\label{lemma-adjoint-functorial-injectives}
Let $\mathcal{A}$ and $\mathcal{B}$ be abelian categories.
Let $u : \mathcal{A} \to \mathcal{B}$ and
$v : \mathcal{B} \to \mathcal{A}$ be additive functors.
Assume
\begin{enumerate}
\item $u$ is right adjoint to $v$,
\item $v$ transforms injective maps into injective maps,
\item $\mathcal{A}$ has enough injectives,
\item $vB = 0$ implies $B = 0$ for any $B \in \Ob(\mathcal{B})$, and
\item $\mathcal{A}$ has functorial injective embeddings.
\end{enumerate}
Then $\mathcal{B}$ has functorial injective embeddings.
\end{lemma}

\begin{proof}
Let $A \mapsto (A \to J(A))$ be a functorial
injective embedding on $\mathcal{A}$. Then
$B \mapsto (B \to uJ(vB))$ is a functorial
injective embedding on $\mathcal{B}$. Compare with the
proof of Lemma \ref{lemma-adjoint-enough-injectives}.
\end{proof}

\begin{lemma}
\label{lemma-partially-defined-adjoint}
Let $\mathcal{A}$ and $\mathcal{B}$ be abelian categories.
Let $u : \mathcal{A} \to \mathcal{B}$ be a functor.
If there exists a subset $\mathcal{P} \subset \Ob(\mathcal{B})$
such that
\begin{enumerate}
\item every object of $\mathcal{B}$ is a quotient of an element
of $\mathcal{P}$, and
\item for every $P \in \mathcal{P}$ there exists an object
$Q$ of $\mathcal{A}$ such that
$\Hom_\mathcal{A}(Q, A) = \Hom_\mathcal{B}(P, u(A))$ functorially
in $A$,
\end{enumerate}
then there exists a left adjoint $v$ of $u$.
\end{lemma}

\begin{proof}
By the Yoneda lemma (Categories, Lemma \ref{categories-lemma-yoneda})
the object $Q$ of $\mathcal{A}$ corresponding to $P$ is defined up to
unique isomorphism by the formula
$\Hom_\mathcal{A}(Q, A) = \Hom_\mathcal{B}(P, u(A))$. Let us write
$Q = v(P)$. Denote $i_P : P \to u(v(P))$ the map corresponding to
$\text{id}_{v(P)}$ in $\Hom_\mathcal{A}(v(P), v(P))$. Functoriality
in (2) implies that the bijection is given by
$$
\Hom_\mathcal{A}(v(P), A) \to \Hom_\mathcal{B}(P, u(A)),\quad
\varphi \mapsto u(\varphi) \circ i_P
$$
For any pair of elements $P_1, P_2 \in \mathcal{P}$ there is a canonical map
$$
\Hom_\mathcal{B}(P_2, P_1)
\to
\Hom_\mathcal{A}(v(P_2), v(P_1)),\quad
\varphi \mapsto v(\varphi)
$$
which is characterized by the rule
$u(v(\varphi)) \circ i_{P_2} = i_{P_1} \circ \varphi$ in
$\Hom_\mathcal{B}(P_2, u(v(P_1)))$.
Note that $\varphi \mapsto v(\varphi)$ is
compatible with composition; this can be seen directly
from the characterization. Hence $P \mapsto v(P)$ is a functor from
the full subcategory of $\mathcal{B}$ whose objects are the elements
of $\mathcal{P}$.

\medskip\noindent
Given an arbitrary object $B$ of $\mathcal{B}$ choose an exact sequence
$$
P_2 \to P_1 \to B \to 0
$$
which is possible by assumption (1). Define $v(B)$ to be the object of
$\mathcal{A}$ fitting into the exact sequence
$$
v(P_2) \to v(P_1) \to v(B) \to 0
$$
Then
\begin{align*}
\Hom_\mathcal{A}(v(B), A)
& =
\Ker(\Hom_\mathcal{A}(v(P_1), A) \to \Hom_\mathcal{A}(v(P_2), A)) \\
& =
\Ker(\Hom_\mathcal{B}(P_1, u(A)) \to \Hom_\mathcal{B}(P_2, u(A))) \\
& =
\Hom_\mathcal{B}(B, u(A))
\end{align*}
Hence we see that we may take $\mathcal{P} = \Ob(\mathcal{B})$, i.e., we
see that $v$ is everywhere defined.
\end{proof}




